% !TeX root = ./thuthesis-example.tex

% 论文基本信息配置

\thusetup{
  %******************************
  % 注意：
  %   1. 配置里面不要出现空行
  %   2. 不需要的配置信息可以删除
  %   3. 建议先阅读文档中所有关于选项的说明
  %******************************
  %
  % 输出格式
  %   选择打印版（print）或用于提交的电子版（electronic），前者会插入空白页以便直接双面打印
  %
  output = electronic,
  %
  % 标题
  %   可使用"\\"命令手动控制换行
  %
  title  = {基于协同引导的多模态生成式推荐方法设计},
  title* = {Design of Collaborative-Guided Multimodal Generative Recommendation Method},
  %
  % 培养单位
  %   填写所属院系的全名
  %
  department = {软件学院},
  %
  % 学科/专业
  %   本科生填写专业名称
  %
  discipline  = {软件工程},
  discipline* = {Software Engineering},
  %
  % 姓名
  %
  author  = {张镕州},
  author* = {To be filled},
  %
  % 指导教师
  %   中文姓名和职称之间以英文逗号","分开
  %
  supervisor  = {张慧, 副研究员},
  supervisor* = {Professor To-be-filled},
  %
  % 日期
  %   使用 ISO 格式；默认为当前时间
  %
  % date = {2026-06-01},
  %
  % 是否在中文封面后的空白页生成书脊（默认 false）
  %
  include-spine = false,
}

% 载入所需的宏包

% 定理类环境宏包
\usepackage{amsthm}

\thusetup{
  math-font = xits,
}

% 表格加脚注
\usepackage{threeparttable}

% 表格中支持跨行
\usepackage{multirow}

% 跨页表格
\usepackage{longtable}

% 算法
\usepackage{algorithm}
\usepackage{algorithmic}

% 量和单位
\usepackage{siunitx}

% 参考文献使用 BibTeX + natbib 宏包
% 本科生参考文献的著录格式
\usepackage[sort]{natbib}
\bibliographystyle{thuthesis-bachelor}

% 定义所有的图片文件在 figures 子目录下
\graphicspath{{figures/}}

% 数学命令
\makeatletter
\newcommand\dif{%  % 微分符号
  \mathop{}\!%
  \ifthu@math@style@TeX
    d%
  \else
    \mathrm{d}%
  \fi
}
\makeatother

% 调整浮动体参数，避免图片单独占据一整页
\renewcommand{\textfraction}{0.15}       % 页面中至少15%是文本
\renewcommand{\topfraction}{0.85}        % 页面顶部最多85%是浮动体
\renewcommand{\bottomfraction}{0.70}     % 页面底部最多70%是浮动体
\renewcommand{\floatpagefraction}{0.66}  % 浮动页至少要有66%被浮动体占据
\setlength{\textfloatsep}{10pt plus 2pt minus 4pt}  % 减小图片与文字的间距
\setlength{\intextsep}{10pt plus 2pt minus 4pt}     % 减小图文混排间距

% hyperref 宏包在最后调用
\usepackage{hyperref}
